% ============================================================
\chapter{Modul 4: Funktionen}
% ============================================================

Ein Roboterprogramm ohne Funktionen ist wie eine Fabrik ohne Maschinen -- alles wird mühsam von Hand erledigt. Funktionen sind wiederverwendbare, benannte Einheiten -- das Fundament strukturierten Codes.

\section{Lektion 11: Eigene Funktionen definieren}

\begin{lstlisting}
def roboter_start():
    print("Systeme pruefen...")
    print("Sensoren initialisiert.")
    print("Motoren bereit.")
    print("Roboter betriebsbereit.")

# Aufruf -- so oft wie nötig:
roboter_start()
\end{lstlisting}

\begin{merksatz}
\texttt{def FUNKTIONSNAME():} definiert die Funktion. Der eingerückte Block darunter ist der Funktionskoerper. Die Funktion läuft \textbf{erst beim Aufruf} -- nicht beim Definieren.
\end{merksatz}

\section{Lektion 12: Parameter}

Parameter machen Funktionen flexibel. Statt einer Funktion pro Geschwindigkeit schreibt man eine Funktion mit Geschwindigkeit als Parameter.

\begin{lstlisting}
def fahre(richtung, geschwindigkeit=0.3, dauer=5):
    weg = geschwindigkeit * dauer
    print(f"Fahre {richtung} mit {geschwindigkeit} m/s fuer {dauer} s")
    print(f"Zurueckgelegter Weg: {weg:.2f} m")

fahre("vorwaerts")               # Standardwerte
fahre("links", 0.15, 3)          # alle Werte
fahre("rueckwaerts", dauer=2)    # gemischt

def drehe(winkel, winkel_geschw=30.0):
    dauer = abs(winkel) / winkel_geschw
    richtung = "links" if winkel > 0 else "rechts"
    print(f"Drehe {abs(winkel)} Grad {richtung} ({dauer:.1f} s)")

drehe(90)
drehe(-45, 45.0)
\end{lstlisting}

\begin{merksatz}
Default-Parameter (mit \texttt{=}) müssen immer \textbf{nach} den Pflichtparametern stehen. \texttt{def fahre(geschwindigkeit=0.3, richtung)} ist ein Fehler. \texttt{def fahre(richtung, geschwindigkeit=0.3)} ist korrekt.
\end{merksatz}

\section{Lektion 13: Rückgabewerte}

\begin{lstlisting}
def berechne_weg(v, t):
    return v * t

def kinematik(v_links, v_rechts, radabstand, dt):
    v = (v_links + v_rechts) / 2
    omega = (v_rechts - v_links) / radabstand
    return v, omega   # zwei Rueckgabewerte als Tupel

# Verwendung:
weg = berechne_weg(0.3, 5.0)
print(f"Weg: {weg} m")

v, omega = kinematik(0.3, 0.35, 0.15, 0.1)
print(f"Lineare Geschwindigkeit: {v:.3f} m/s")
print(f"Winkelgeschwindigkeit:   {omega:.3f} rad/s")
\end{lstlisting}

\begin{praxis}
Die \texttt{kinematik()}-Funktion ist der Kern jeder Differentialantrieb-Steuerung. Mit den Radgeschwindigkeiten links und rechts sowie dem Radabstand berechnet sie die Bewegung des Roboters. Diese Formel wirst du in ROS2 wiedersehen.
\end{praxis}

\section{Lektion 14: Programme strukturieren}

\begin{lstlisting}
# Vollstaendiges, strukturiertes Roboterprogramm

def initialisierung():
    print("System initialisiert.")
    return True

def lese_sensor(sensor_id):
    # Simulierter Sensorwert
    import random
    return round(random.uniform(0.1, 2.0), 2)

def entscheide(abstand):
    if abstand < 0.20:
        return "stopp"
    elif abstand < 0.50:
        return "langsam"
    else:
        return "normal"

def fuehre_aus(befehl):
    geschwindigkeiten = {"stopp": 0.0, "langsam": 0.1, "normal": 0.3}
    v = geschwindigkeiten.get(befehl, 0.0)
    print(f"Befehl: {befehl:8} | Geschwindigkeit: {v} m/s")

def hauptschleife(iterationen=10):
    for i in range(iterationen):
        abstand = lese_sensor("ultraschall")
        befehl  = entscheide(abstand)
        fuehre_aus(befehl)

# Hauptprogramm
if initialisierung():
    hauptschleife(10)
\end{lstlisting}

\begin{merksatz}
\textbf{Eine Funktion, eine Aufgabe.} \texttt{lese\_sensor()} liest, \texttt{entscheide()} entscheidet, \texttt{fuehre\_aus()} handelt. Diese Trennung nennt man \textit{Separation of Concerns} -- sie macht Code testbar, wartbar und erweiterbar.
\end{merksatz}

\begin{ros}
In ROS2-Nodes sieht die Struktur genau so aus: \texttt{\_\_init\_\_()} initialisiert, \texttt{timer\_callback()} ist die Hauptschleife, \texttt{sensor\_callback()} verarbeitet Eingaben. Jede Funktion hat genau eine Aufgabe.
\end{ros}

\begin{uebung}[title=Projektaufgabe Modul 4: Roboter-Bibliothek]
Schreibe ein strukturiertes Programm mit mindestens 6 Funktionen:
\texttt{starte()}, \texttt{stoppe()}, \texttt{fahre\_vorwaerts(v, t)}, \texttt{drehe(winkel)}, \texttt{lese\_abstand()}, \texttt{fahre\_route(punkte)}.
\texttt{fahre\_route()} nimmt eine Liste von (Strecke, Winkel)-Tupeln und ruft die anderen Funktionen auf.
\end{uebung}
